\\newpage
\\section{模型的建立与求解}

\\subsection{数据预处理}
\\begin{enumerate}[label=(\\arabic*), itemindent=0pt, labelindent=\\parindent, labelwidth=2em, labelsep=5pt, leftmargin=*]
    \\item[\ 1$] \quad 缺失值检测与线性插值补全
    \\item[\ 2$] \quad 异常值检测（3\\sigma准则）与Winsorize处理
    \\item[\ 3$] \quad Min-Max标准化：^* = \\frac{x - x_{\\min}}{x_{\\max} - x_{\\min}}$
\\end{enumerate}

\\subsection{问题一模型：MNL选择模型}
假设车主选择停车场：
U_i = \\alpha - \\beta c_i + \\epsilon_i
其中，$\\epsilon_i，服从Gumbel分布。根据离散选择理论，选择概率为：
P_i = \\frac{e^{U_i}}{\\sum_j e^{U_j}} = \\frac{e^{\\alpha - \\beta c_i}}{\\sum_j e^{\\alpha - \\beta c_j}}

\\subsection{问题二模型：Stackelberg博弈}
\\subsubsection{领导者问题（运营商）}
\\max_p \\pi_1 = \\sum_i (c_i - cost_i) \\cdot D_i(p)
\\text{s.t.} \\quad c_i^{\\min} \\leq c_i \\leq c_i^{\\max}

\\subsubsection{追随者问题（车主）}
q_i^* = \\arg\\max_q \\pi_2 = \\sum_i U_i(q_i)
\\text{s.t.} \\quad \\sum_i q_i = Q_{total}

\\subsubsection{Stackelberg均衡}
均衡满足：
p^* = \\arg\\max_p \\pi_1(p, q^*(p))
q^* \\in \\arg\\max_q \\pi_2(p^*, q)

\\subsection{模型求解}
采用KKT条件求解均衡解。通过数值搜索方法，得到最优定价策略。

\\subsection{求解结果}
\\begin{table}[H]
\\centering
\\caption{最优定价策略}
\\begin{tabular}{c c c}
\\toprule
停车场 & 最优价格(元/小时) & 预期需求量 \\
\\midrule
A & 15.2 & 180 \\
B & 12.8 & 220 \\
C & 18.5 & 150 \\
\\bottomrule
\\end{tabular}
\\end{table}
